\begin{textsample}{POS Dim 4 – vem\_ed\_18 – Score 5.00 – vem\_ed\_18/t081.txt}  \label{ex:f4_pos_045}
AOS LEITORES AOS LEITORES No dia 8 de dezembro de 2022, a equipe de PCIs / Cesu, com o apoio das equipes EDI e Formação Continuada da Cesu, realizou o primeiro Simpósio Internacional de Intercâmbios Virtuais (IVES Cesu 2022).
O evento, realizado nos períodos da manhã e tarde, reuniu 5 \textbf{especialistas} internacionais, 9 professores do Centro Paula Souza e um público total de 148 participantes.
A síntese das apresentações está em VEm 15 (https://cesu.cps.sp.gov.br/wp-content/uploads/2022/12/VEm-15.pdf).
A partir daquela experiência, a equipe de PCIs / Cesu entendeu que o IVES Cesu 2023 deveria ser realizado no primeiro semestre, com foco na \textbf{divulgação} das experiências institucionais e em temas de \textbf{destaque} no momento, contando com a presença de \textbf{especialistas} internacionais em Intercâmbios Virtuais.
Esta edição de VEm é inteiramente dedicada à cobertura do IVES Cesu 2023, que ocorreu em 25 de \textbf{maio}, das 14 às 17h.
Participaram do evento 190 pessoas, sendo 93 de Fatecs, 42 de Etecs, 7 do Centro Paula Souza, 21 de outras instituições educacionais brasileiras e 31 de universidades estrangeiras.
A maior parte (151) é de brasileiros, porém houve participantes também dos seguintes países: México (12), EUA (10), Chile (4), Espanha (4), Colômbia (3), Equador (2), Canadá (2), França (1) e Quênia (1).
As páginas a seguir resumem as principais discussões trazidas pelos coordenadores de Intercâmbios Virtuais da DePaul University (EUA), DUOC UC (Chile), Universidad Católica Andrés Bello (Venezuela), University of Minnesota System e SUNY COIL Center (EUA).
Boa leitura!
EXPEDIENTE Expediente CPS Diretora-Superintendente: Laura Laganá Vice-Diretora-Superintendente: Emilena Lorenzon Bianco Chefe de Gabinete: Armando Natal Maurício Expediente Cesu Coordenador Técnico: Rafael Ferreira Alves Diretor Acadêmico-Pedagógico: André Luiz Braun Galvão Gestão Educacional: Willian Marcos Muniz Menezes Departamento Administrativo: Silvia Pereira Abranches EDI - Estruturação e Desenvolvimento Instrucional: Thais Lari Braga Cilli Expediente Línguas e Projetos Colaborativos Internacionais - Cesu Coordenação de Línguas e Projetos Internacionais: Mariane Teixeira Coordenação de Projetos Colaborativos Internacionais: Osvaldo Succi Junior Acompanhamento pedagógico PCI: Ana Carolina Freschi, Neusa Haruka Gritti e Regiane Moreira Expediente VEm Corpo editorial: Ana Carolina Freschi, Mariane Teixeira, Neusa Haruka Gritti, Osvaldo Succi Junior e Regiane Moreira Jornalista responsável e Comunicação: Patrícia Patricio - MTb 25.131 Editoração e diagramação: Fábio Gomes da Silva VEm: Virtual Exchange Medium é um informativo com publicação bimestral da Cesu / CEETEPS: Rua dos Andradas, 140 - Santa Efigênia - 01208-000 - São Paulo - SP

% matched lemmas: destaque, divulgação, especialista, maio
\end{textsample}
